\section{Method}

\subsection{Problem Setting}

We study three forms of temporal reuse on frozen video-VLM stacks:

\begin{itemize}[leftmargin=1.5em]
\item \textbf{First-pass visual pruning:} skip or reduce vision-tower work
before the rest of the model pays for it.
\item \textbf{After-ingest follow-up reuse:} keep the expensive prefix from the
first query and ask the next question against the same video.
\item \textbf{Routing under a fixed dense backend:} substitute cached visual
state to test whether answer quality survives, even when the benchmark path has
not yet been rewritten into a sparse wall-clock implementation.
\end{itemize}

These are related, but they are not interchangeable. They are kept
analytically separate because they measure different kinds of reuse.

\subsection{Temporal Planner}

For adjacent frames, the routing path computes a per-block statistic on the RGB
difference plane. The implementation supports mean difference, max-absolute
difference, changed-pixel fraction, and top-$k$ mean over the block. In the
best current clean-tree Qwen policy, the statistic is \texttt{max\_abs}. Fixed
thresholds partition blocks into \emph{static}, \emph{shifted}, and
\emph{novel} regions. These names are historical planner labels, not literal
codec semantics: today they mean low, medium, and high pixel-domain change.

The block geometry is derived from the model's vision configuration rather than
hard-coded. On Qwen2.5-VL, planner decisions line up with the merged-token
spatial grid so cached features can be substituted back into the dense path at
the same geometry the benchmark runner already understands.

\subsection{Bounded Staleness}

Reuse without age control drifts. The planner therefore carries a
bounded-staleness counter: even a block that continues to look reusable must be
refreshed after its age exceeds a configured limit. The clean-tree policy used
in the paper sets that limit to four frame steps. That choice is not cosmetic.
It is the main reason the planner survives TOMATO, where temporally concentrated
evidence punishes policies that keep trusting a stale region simply because the
local pixel difference remains small.

We summarize routing budget by converting block-level reuse into an effective
fresh-frame count,
\[
f_{\mathrm{eff}} = 1 + (N - 1)(1 - r_{\mathrm{reuse}}),
\]
where $r_{\mathrm{reuse}}$ is the mean active-region reuse ratio. This gives
the Qwen routing results a common x-axis with dense-$N$ baselines: how much
fresh visual budget the policy effectively spent.

\subsection{Vision-Tower Pruning and Persistent-KV}

The other two mechanisms operate at different points in the stack.

The Gemma path patches the vision tower after an internal layer, keeps only a
fraction of the tokens, and scatters the result back so downstream geometry is
unchanged. The resulting first-pass speedup is constrained by how much of total
wall-clock the vision tower owns. That is why the paper reports the
scatter-back ceiling
\[
\mathrm{E2E} \le \frac{1}{1 - V_{\mathrm{share}} V_{\mathrm{red}}},
\]
where $V_{\mathrm{share}}$ is the dense vision share and $V_{\mathrm{red}}$ is
the observed reduction in vision time.

The persistent-KV path answers a different question. Here the user has already
paid for the first query on a video. The next question about that same video
reuses the prompt cache and only appends a short textual suffix. The result can
be dramatically faster, but it is explicitly an \emph{after-ingest} claim. It
is not a claim about faster first-pass understanding of fresh videos.

\subsection{What Is Actually Skipped}

The paper keeps semantic substitution and wall-clock skipping separate. The
Qwen routing runner can replay or substitute cached features to test whether the
answer survives. That is enough to study the quality--compute frontier. It is
not enough to claim an end-to-end speedup until the backend itself skips timed
work. The Gemma and persistent-KV lanes do change measured wall-clock and
therefore support speedup claims. The Qwen routing lane does not yet. That
choice is deliberate.

\begin{sectiontodo}{Add the overview figure and final algorithm box}
The prose now matches the paper's current scope, but the final manuscript still
needs a compact visual that makes three distinctions obvious: first-pass
pruning, after-ingest prompt reuse, and Qwen routing as mechanism validation
rather than measured sparse execution.

\textcolor{todoBorder}{Source pointers: \path{src/codec_through/temporal.py},
\path{src/codec_through/novelty_pruning.py},
\path{src/codec_through/pruned_vision_tower.py},
\path{scripts/run_kv_cache_session.py},
\path{paper/claim-matrix.md}.}
\end{sectiontodo}
